\section*{Chapitre \ref{ch:g6:classification-kinship} --- Classification et parenté}

\begin{solution}{exo:g6:classification-kinship:1}
Un caractère qu'un organisme possède, énoncé de façon vérifiable~:
possède une \omterm{def:g4:classifying-animals:vertebrate}{colonne vertébrale}~; possède quatre \omterm{def:g1:my-body:parts}{membres}~; possède
des poils et allaite ses petits (ou~: possède des \omterm{def:g1:animals-around-us:covering}{plumes}).
\end{solution}

\begin{solution}{exo:g6:classification-kinship:2}
L'ensemble des possesseurs de chaque \omterm{def:g6:classification-kinship:attribute}{attribut} se range tout entier
dans celui d'un autre~: tous les porteurs de poils ont quatre
\omterm{def:g1:my-body:parts}{membres}, tous les \omterm{def:g1:animals-around-us:animal}{animaux} à quatre \omterm{def:g1:my-body:parts}{membres} ont une \omterm{def:g4:classifying-animals:vertebrate}{colonne
vertébrale} --- jamais l'inverse, jamais de croisement. Des boîtes
dans des boîtes, sans aucun cas à cheval.
\end{solution}

\begin{solution}{exo:g6:classification-kinship:3}
Parce qu'elles les ont hérités d'un ancêtre commun~: un \omterm{def:g6:classification-kinship:attribute}{attribut}
descend le long d'une lignée familiale, si bien que ses possesseurs
sont les descendants de cet ancêtre --- une famille.
\end{solution}

\begin{solution}{exo:g6:classification-kinship:4}
La truite se tient au-delà du seul point de la \omterm{def:g4:classifying-animals:vertebrate}{colonne vertébrale}.
La baleine se tient au-delà de la colonne, des quatre \omterm{def:g1:my-body:parts}{membres} et des
poils et \omterm{def:g2:animals-grow-up:born}{lait} --- les trois.
\end{solution}

\begin{solution}{exo:g6:classification-kinship:5}
La baleine et le chat~: ils partagent la fourche des poils et du
\omterm{def:g2:animals-grow-up:born}{lait}, la plus récente qui soit dessinée~; la truite s'est séparée à
la plus ancienne.
\end{solution}

\begin{solution}{exo:g6:classification-kinship:6}
C'est deux fois la même information~: une boîte est une \omterm{def:g1:plants-around-us:tree}{branche} avec
tous ses rameaux --- tout ce qui se trouve au-delà d'une fourche
hérite de son \omterm{def:g6:classification-kinship:attribute}{attribut} et se range dans sa boîte.
\end{solution}

\begin{solution}{exo:g6:classification-kinship:7}
Les poils et le \omterm{def:g2:animals-grow-up:born}{lait} sont hérités le long d'une seule lignée
familiale --- ils marquent l'ascendance. «~A la forme d'un
poisson~» peut être imprimé sur des étrangers par un même mode de
vie aquatique, comme le montre la baleine~: forme empruntée, famille
ailleurs. La \omterm{def:g4:classifying-animals:classification}{classification} trie des familles~: elle se fie donc à
l'héritage et se méfie de l'habit emprunté.
\end{solution}

\begin{solution}{exo:g6:classification-kinship:8}
Truite~: colonne seulement. Grenouille~: colonne, quatre \omterm{def:g1:my-body:parts}{membres}.
Chat~: colonne, quatre \omterm{def:g1:my-body:parts}{membres}, poils et \omterm{def:g2:animals-grow-up:born}{lait}. Merle~: colonne,
quatre \omterm{def:g1:my-body:parts}{membres}, \omterm{def:g1:animals-around-us:covering}{plumes}. L'emboîtement se voit dans les croix~: les
croix de chaque ligne comprennent la colonne~; toute ligne qui coche
un revêtement coche aussi les quatre \omterm{def:g1:my-body:parts}{membres}.
\end{solution}

\begin{solution}{exo:g6:classification-kinship:9}
Il franchit la \omterm{def:g4:classifying-animals:vertebrate}{colonne vertébrale} et les quatre \omterm{def:g1:my-body:parts}{membres}, et part sur
sa propre \omterm{def:g1:plants-around-us:tree}{branche} avec sa \omterm{def:g1:the-five-senses:sense}{peau} sèche et écailleuse --- à côté, et
non au-delà, des fourches des poils et des \omterm{def:g1:animals-around-us:covering}{plumes}.
\end{solution}

\begin{solution}{exo:g6:classification-kinship:10}
L'allaitement marque la \omterm{prop:g6:classification-kinship:kinship}{parenté}~: les poils et le \omterm{def:g2:animals-grow-up:born}{lait} sont
l'héritage de la famille des \omterm{prop:g3:sorting-living-things:groups}{mammifères}, si bien que la
chauve-souris se \omterm{def:g1:plants-around-us:tree}{branche} avec le chat. Le vol est emprunté au mode
de vie --- les ailes sont apparues séparément sur la \omterm{def:g1:plants-around-us:tree}{branche} des
\omterm{prop:g3:sorting-living-things:groups}{oiseaux} et sur celle des chauves-souris, comme le corps fuselé chez
la baleine et la truite.
\end{solution}

\begin{solution}{exo:g6:classification-kinship:11}
Il lui appartient. L'ornithorynque descend de l'ancêtre commun des
\omterm{prop:g3:sorting-living-things:groups}{mammifères} et porte les héritages de la famille (poils, \omterm{def:g2:animals-grow-up:born}{lait})~; son
\omterm{def:g2:animals-grow-up:born}{œuf} est un caractère ancien que sa \omterm{def:g1:plants-around-us:tree}{branche} précoce a conservé alors
que les autres l'ont perdu. La famille est l'ancêtre et \emph{tous}
ses descendants --- bizarreries comprises~; le paquet de caractères
(\cref{exo:g4:classifying-animals:11}) l'avait déjà admis.
\end{solution}

\begin{solution}{exo:g6:classification-kinship:12}
Si chaque paire de \omterm{def:g1:plants-around-us:tree}{branches} se rejoint, alors les \omterm{def:g5:biodiversity:species}{espèces} doivent
changer au fil des générations sur de longues durées --- de nouveaux
\omterm{def:g6:classification-kinship:attribute}{attributs} apparaissant et étant hérités --- pour qu'une lignée
ancestrale ait pu se ramifier en truite, en chat et en pâquerette.
Implication en une phrase~: tous les \omterm{def:g6:exploring-our-environment:components}{êtres vivants} descendent, avec
changement, d'ancêtres communs. Preuves à exiger~: des traces du
changement --- des formes intermédiaires, des restes conservés dans
des roches anciennes, et des airs de famille placés exactement là où
l'\omterm{def:g6:classification-kinship:tree}{arbre} le prédit (l'affaire d'une année ultérieure).
\end{solution}

\begin{solution}{pb:g6:classification-kinship:1}
\textbf{1.} Truite~: \omterm{def:g4:classifying-animals:vertebrate}{colonne vertébrale}. Grenouille~: colonne,
quatre \omterm{def:g1:my-body:parts}{membres}. Lézard~: colonne, quatre \omterm{def:g1:my-body:parts}{membres}, \omterm{def:g1:the-five-senses:sense}{peau} sèche et
écailleuse. Pigeon~: colonne, quatre \omterm{def:g1:my-body:parts}{membres}, \omterm{def:g1:animals-around-us:covering}{plumes}. Souris~:
colonne, quatre \omterm{def:g1:my-body:parts}{membres}, poils et \omterm{def:g2:animals-grow-up:born}{lait}. Dauphin~: colonne, quatre
\omterm{def:g1:my-body:parts}{membres}, poils et \omterm{def:g2:animals-grow-up:born}{lait} (le témoignage de l'étiquette sur le \omterm{def:g2:animals-grow-up:born}{lait}~;
ce sont les caractères qui décident, et le \omterm{def:g2:animals-grow-up:born}{lait} est le caractère
décisif).

\textbf{2.} La \omterm{def:g4:classifying-animals:vertebrate}{colonne vertébrale} --- tous les six sont des
\omterm{def:g4:classifying-animals:vertebrate}{vertébrés}. Le partage le plus large marque l'ancêtre commun le plus
ancien~: à cette profondeur, les six ne forment qu'une famille.

\textbf{3.} La grenouille, le lézard, le pigeon, la souris, le
dauphin. Les ailes du pigeon et les nageoires du dauphin sont quatre
\omterm{def:g1:my-body:parts}{membres} en habit de travail~: les mêmes \omterm{def:g1:my-body:parts}{membres} hérités, façonnés
par des vies différentes --- l'\omterm{def:g6:classification-kinship:tree}{arbre} compte l'héritage, pas l'habit.

\textbf{4.} Parce qu'une adresse partagée n'est pas une famille
partagée~: «~vit dans l'eau~» est un milieu, empruntable par des
étrangers, et non un \omterm{def:g6:classification-kinship:attribute}{attribut} de l'ascendance du corps --- la
vieille règle de \cref{def:g3:sorting-living-things:criterion},
munie maintenant de sa raison.

\textbf{5.} Tronc $\rightarrow$ fourche de la colonne (la truite
part après elle) $\rightarrow$ fourche des quatre \omterm{def:g1:my-body:parts}{membres}
$\rightarrow$ trois \omterm{def:g1:plants-around-us:tree}{branches} de revêtement~: \omterm{def:g1:animals-around-us:covering}{écailles} sèches
(lézard), \omterm{def:g1:animals-around-us:covering}{plumes} (pigeon), poils et \omterm{def:g2:animals-grow-up:born}{lait} (souris et dauphin
ensemble sur la fourche la plus récente)~; la grenouille part après
les quatre \omterm{def:g1:my-body:parts}{membres}, avant tout revêtement.

\textbf{6.} Le cousin de tiroir le plus proche de la souris est le
dauphin (ils partagent la fourche des poils et du \omterm{def:g2:animals-grow-up:born}{lait}). Celui de la
truite, c'est tout le monde à égalité --- elle s'est séparée à la
première fourche, si bien que les cinq autres sont des cousins
également éloignés.

\textbf{7.} Un amphibien --- quatre \omterm{def:g1:my-body:parts}{membres}, \omterm{def:g1:the-five-senses:sense}{peau} nue et humide,
aucune fourche de revêtement à elle~; sa \omterm{def:g1:plants-around-us:tree}{branche} part entre la
fourche des quatre \omterm{def:g1:my-body:parts}{membres} et les revêtements.

\textbf{8.} Au-delà de la colonne et des quatre \omterm{def:g1:my-body:parts}{membres}, sur la
\omterm{def:g1:plants-around-us:tree}{branche} des \omterm{def:g1:animals-around-us:covering}{écailles} sèches~: un reptile --- classé sans être vu,
parce que ce sont les \omterm{def:g6:classification-kinship:attribute}{attributs}, non les étiquettes, qui décident.

\textbf{9.} Le dauphin et la souris se tiennent plus près~: les
poils et le \omterm{def:g2:animals-grow-up:born}{lait} sont une fourche héritée. Le corps fuselé est
l'uniforme de l'eau, distribué pareillement à des nageurs sans
\omterm{prop:g6:classification-kinship:kinship}{parenté} (\cref{exo:g4:adapted-to-their-world:11})~; compté, il
classerait un habit emprunté comme une famille et ruinerait
l'emboîtement.

\textbf{10.} Une unique \omterm{def:g5:biodiversity:species}{espèce} ancestrale --- un ancien \omterm{def:g4:classifying-animals:vertebrate}{vertébré} à
quatre \omterm{def:g1:my-body:parts}{membres} --- dont la grenouille, le lézard, le pigeon, la
souris et le dauphin descendent tous, chacun portant ses quatre
\omterm{def:g1:my-body:parts}{membres} en habit différent.

\textbf{11.} Que tous les six --- et la classe qui les dessine ---
appartiennent à une seule famille ramifiée, descendue avec
changement d'ancêtres communs, les humains se trouvant sur la
\omterm{def:g1:plants-around-us:tree}{branche} des poils et du \omterm{def:g2:animals-grow-up:born}{lait} parmi les autres.

\textbf{12.} L'emboîtement est un fait du monde vivant --- les
\omterm{def:g6:classification-kinship:attribute}{attributs} se répartissent réellement en boîtes emboîtées, ce
qu'aucun archiviste n'a décrété. L'\omterm{def:g6:classification-kinship:tree}{arbre} explique ce fait~:
l'hérédité à partir d'ancêtres communs \emph{doit} produire des
\omterm{def:g6:classification-kinship:attribute}{attributs} emboîtés, alors qu'un simple classement laisserait
l'obligeante propreté du monde à l'état de miracle.
\end{solution}
